\section{Quantifying Representational Harms}
\label{sec:formalization}

% \begin{figure*}[h]
% \begin{subfigure}[b]{0.24\textwidth}
% \includegraphics[width=\textwidth,trim=0cm 0cm 0cm 0cm,clip=true]{./images/conceptnet_sentiment_negative.pdf}
% \end{subfigure}
% \begin{subfigure}[b]{0.24\textwidth}
% \includegraphics[width=\textwidth,trim=0cm 0cm 0cm 0cm,clip=true]{./images/conceptnet_sentiment_positive.pdf}
% \end{subfigure}
% \begin{subfigure}[b]{0.24\textwidth}
% \includegraphics[width=\textwidth,trim=0cm 0cm 0cm 0cm,clip=true]{./images/conceptnet_regard_negative.pdf}
% \end{subfigure}
% \begin{subfigure}[b]{0.24\textwidth}
% \includegraphics[width=\textwidth,trim=0cm 0cm 0cm 0cm,clip=true]{./images/conceptnet_regard_positive.pdf}
% \end{subfigure}

% \caption{Negative sentiment and regard results from ConceptNet.}
% \label{fig:Bias_results}
% \end{figure*}

%define bias for a set of statement each targeting some subject
%Section 2: 
%Para1: talk about bias in high level, formalize a set of sentences with notations, say we use sentiment and regard as approximations for prejudice/favoritism and introduce two subtypes: Overgeneralization and Disparity. 
%Para2: Overgeneralization formulation: average percentages of approximated prejudice/favoritism. 
%Para3: Disparity formulation: representation difference (frequency) and prejudice/favoritism difference
%define notation for “target subject” mentioned in the sentence, or whatever name that is appropriate; and use that for later formulation, etc


%This section introduces our definition and measurements of bias in commonsense knowledge resources.
%Generally, bias is the presence of prejudice or favoritism toward an individual or a demographic group \cite{mehrabi2019survey}. 
%We focus on biases towards demographic groups of multiple social categories. For example, ``\emph{Chinese}'' and ``\emph{American}'' are two groups in the ``\emph{race}'' category.


\subsection{Representational Harms}

% \para{Representational Harms}
Representational harms occur ``\emph{when systems reinforce the subordination of some groups along the lines of identity}'' and can be further categorized into stereotyping, recognition, denigration, and under-representation~\cite{barocas2017problem, crawford2017trouble,  blodgett-etal-2020-language}. 
% It has been argued that compared to allocative harms that occur ``\emph{when a system allocates or withholds certain groups an opportunity or a resource}'', representational harms have deeper long-term effects and are more challenging to quantify~\cite{crawford2017trouble, blodgett-etal-2020-language}.
This work aims to formalize representational harms specifically for \emph{a set of statements} about some target groups~\cite{nadeem2020stereoset}, e.g ``\emph{lawyer is related to dishonest}'' is a statement about the group ``\emph{lawyer.}'' 

When measuring such harms, we consider the core concept of \emph{polarized perceptions}: non-neutral views that can take the form of either prejudice that expresses negative views~\footnote{The word \emph{prejudice} is defined as ``\emph{preconceived (usually unfavorable) evaluation of another person}''~\cite{lindzey1968handbook}. We focus on only unfavorable prejudice and use favoritism for favorable evaluation.} or favoritism that expresses positive views toward a certain target perceived in the statement~\cite{10.1145/3457607}.
%formalize two specific types of \emph{representational harms} that are defined as ``when a system (e.g., a search engine) represents some social groups in a less favorable light than others, demeans them, or fails to recognize their existence altogether''~\cite{ blodgett-etal-2020-language}: \emph{overgeneralization} of polarized perceptions toward each target: $Bias_O(\mathbb{S}, \mathbb{T})$ and \emph{disparity} in representation and overgeneralization across targets: $Bias_D(\mathbb{S}, \mathbb{T})$.
%We define bias specifically for a set of statements each targeting some subject. 
% \para{Notations}

More formally, let $\mathbb{S} = \left \{  s_1, s_2, ..., s_n\right \}$ indicate the set of $n$ natural language statements $s_i$ and let $\mathbb{T} = \left \{  t_1, t_2, ..., t_m,\right \}$ indicate the set of $m$ targets such that each $t_j$ has appeared at least once in $\mathbb{S}$. Each statement $s_i$ contains a target $t_j$. We use $s_i^{+/-}(t_j)$ to indicate that the statement expresses a positive or negative perception toward the target $t_j$.
We are interested in quantifying the representational harms toward $\mathbb{T}$ in this set $\mathbb{S}$.


\subsection{Two Types of Harms}
To adapt the definition of representational harms to a sentence set, we define two sub-types of harms, \emph{intra-target overgeneralization} and \emph{inter-target disparity}, aiming to cover different categories of representational harms~\cite{barocas2017problem, crawford2017trouble}. We consider \emph{overgeneralization} that directly examines whether targets such as ``\emph{lawyer}'' or ``\emph{lady}'' are perceived positively or negatively in the statements (examples in Table~\ref{motivation_table}), covering categories including stereotyping, denigration, and favoritism. Then we consider \emph{disparity} across different targets in representation (do some targets have fewer associated statements and lower coverage) and polarized perceptions (whether some targets are more positively or negatively perceived).
%We would like to study whether triples related to specific groups of people contain prejudices in favor of or against the groups.
%which is harmful in their own rights due to the close association between language usage and stereotyping in human cognitives~\cite{beukeboom2019stereotypes}.


\para{Intra-target Overgeneralization}
The ideal sentence set depicting a target group such as ``\emph{lawyer}'' or ``\emph{lady}'' should have neither \emph{favoritism} nor \emph{prejudice} toward the target group.
Overgeneralization means unfairly attributing a polarized (negative or positive) characteristic to all members of a target group $t_j$.
In the context of sentence sets $\mathbb{S}$, we define overgeneralization to be polarized sentences in the set regarding the targets. When the statement $s_i$ contains negatively-polarized views toward a target group such as ``\emph{lawyers are dishonest,}'' it demonstrates \textit{prejudice} toward lawyers. It is an overgeneralized statement because it implies that \emph{all} lawyers are not honest. The same logic applies to positively-polarized statements such as ``\emph{British people are brilliant}'' that constitute favoritism.
 
 % More discussion on why these two are effective measures
 
 %in Section~\ref{human_eval}.
 %\xiang{I think so far most of this paragraph (the notations and intro about polarization) are generally for the representational harms; not just for overgeneralization; we should re-arrange it to put things to the right place..}
 
% \begin{table}[]
% \centering
% \scalebox{0.93}{
% \begin{tabular}{ p{3.5cm} c }
%  \toprule
% \textbf{}&\textbf{Agreement with Human}\\
%  \midrule
%  \parbox[t]{2mm}{\multirow{1}{*}{\shortstack[l]{GenericsKB Regard}}}
% &60.9\%\\[0.5pt]

%  \parbox[t]{2mm}{\multirow{1}{*}{\shortstack[l]{GenericsKB Sentiment}}}
% &70.3\%\\[0.5pt]

%  \parbox[t]{2mm}{\multirow{1}{*}{\shortstack[l]{ConceptNet Regard}}}
% &75.4\%\\[0.5pt]

%  \parbox[t]{2mm}{\multirow{1}{*}{\shortstack[l]{ConceptNet Sentiment}}}
% &83.1\%\\[0.5pt]
% %  \midrule

% %  \midrule
%  \bottomrule
% \end{tabular}
% }
% \caption{Agreement of sentiment and regard labels with human annotators in terms of accuracy. \xiang{can rotate to put into 3 rows (1st row dataset; 2nd row: sentiment, regard; 3rd row: numbers).}}
% \label{human_vs_re_sent_generics}
% \end{table}




 
%(CAN PUT BACK IF NEEDED) We emphasize that sentiment and regard are far from perfect measures for polarizations in perceptions (favoritism and prejudice), for example, ``\emph{police are capable of arresting people}'' might be classified as non-neutral by a sentiment classifier, but it would not be counted as a polarized perception according to our definition of representational harms since it is a description of the job's duty. Furthermore, our assumption that an ideal set would be \emph{neutral} (no favoritism or prejudice toward any group) should not be conflated with sentiment or regard neutrality. However, these two measures are suitable for our case since they both capture polarity of sentences which can give us approximations for favoritism and prejudice towards a target group. 
%Prior work~\cite{sheng2019woman} has demonstrated that sentiment and regard are effective measures of bias (polarized views towards a target group). Although this is still an active area of research, for now, these are promising proxies that many work in ethical NLP also have used to measure bias (e.g.~\citet{sheng2019woman, li2020unqover, brown2020language, sheng2020towards,dhamala2021bold}). We also conduct additional crowdsourcing experiments that validate this claim in Section~\ref{human_eval}. 
% We use these measures to quantify the overgeneralization score for every target in $\mathbb{S}$, and we would ideally need sentiment and regard to label all $s_i$ in $\mathbb{S}$ as neutral sentences to achieve perfect fairness for every group. %We collect all sentences describing a particular target $t_j$ and measure overgeneralization for this group in the sentence set $\mathbb{S}$.
%To measure polarization in social perceptions towards targets in $\mathbb{T}$ in the set $\mathbb{S}$, we first consider two approximations: 1. The sentiment polarity of triples with positive for favoritism and negative for prejudice. 2. The \textit{Regard} measure which is specifically designed to capture the language polarity towards social perceptions of a demographics \cite{sheng2019woman}. 

To measure the polarization in a sentence $s_i$, we use two approximations: sentiment polarity and regard~\cite{sheng2019woman}. We define an intra-target overgeneralization score using these measures for every target in $\mathbb{S}$. Formally, for each target $t_j$, we collect all statements in $\mathbb{S}$ that contain target $t_j$ to form a target-specific set $\mathbb{S}_{t_j}$, we then apply the sentiment and regard classifiers individually on all statements in $\mathbb{S}_{t_j}$ and report the average percentages of statements that are classified as positive or negative sentiment or regard labels for the target $t_j$. The negative and positive overgeneralization bias (prejudice and favoritism) for $t_j$ is quantified as:
\begin{equation} \label{eq:1}
     O^{-}(\mathbb{S}, t_j) = |\mathbb{S}^{-}_{t_j}||\mathbb{S}_{t_j}|^{-1} \times 100,
    %  \vspace{-1em}
\end{equation}
\begin{equation} \label{eq:2}
     O^{+}(\mathbb{S}, t_j) = |\mathbb{S}^{+}_{t_j}||\mathbb{S}_{t_j}|^{-1} \times 100,
\end{equation}
where $|\mathbb{S}^{+/-}_{t_j}|$ is the number of statements with positive/negative polarization measured by sentiment or regard for the target $t_j$, i.e. $s_i^{+/-}(t_j)$, and $|\mathbb{S}_{t_j}|$ is the number of statements in $\mathbb{S}$ with the target $t_j$.

%During this labeling process using sentiment and regard classifiers, we mask the demographic information from the sentences to avoid biases in sentiment and regard classifiers to affect our analysis. For sentiment analysis, we utilize the VADER sentiment analysis tool \cite{gilbert2014vader} which is a rule based sentiment analyzer and use the BERT model for the regard classifier from \cite{sheng2019woman}. 

%To quantify \emph{overgeneralization}, after converting triples $T$ in the form of $(s,r,o)$ to natural language sentences $S$, we apply sentiment and regard classifiers and report the \emph{percentages of sentences} $s \in S$ that have either positive or negative sentiment (or regard) scores for different target demographic groups $g \in G$ as follows:$\text{Percent Positive Regard}_{g \in G} = \frac{|S^{r+}_g|}{|S_g|} \times 100$ and $\text{Percent Negative Regard}_{g \in G} = \frac{|S^{r-}_g|}{|S_g|} \times 100$, where $|S^{r+/-}_g|$ represents the number of sentences with positive/negative regard labels for the demographic group $g$. Similarly, for sentiment scores, we use $|S^{s+/-}_g|$ representing the number of sentences with positive/negative sentiment labels, and $|S_g|$ the total number of sentences associated with group $g$. For instance, the triple $(\text{American, IsA, great\_person})$, representing the ``\emph{American}'' demographic group $(g=American)$, will be converted to ``\emph{American is a great person}'' sentence on which sentiment and regard classifiers will be applied. Notice in the example provided, ``\emph{American}'' is the demographic group $g$ which belongs to the ``\emph{race}'' category.

% \[
% \text{Percent Negative Regard}_{g \in G} = \frac{|S^{r-}_g|}{|S_g|} \times 100
% \]

% \[
% \text{Percent Positive Sentiment}_{g \in G} = \frac{|S^{s+}_g|}{|S_g|} \times 100
% \]

% \[
% \text{Percent Negative Sentiment}_{g \in G} = \frac{|S^{s-}_g|}{|S_g|} \times 100
% \]


% \begin{figure*}[h]
% % \centering
% \begin{subfigure}[b]{0.33\textwidth}
% \includegraphics[width=\textwidth,trim=7.3cm 0cm 7.3cm 0cm,clip=true]{./images/new_ex_scatter.pdf}
% \end{subfigure}
% \begin{subfigure}[b]{0.33\textwidth}
% \includegraphics[width=\textwidth,trim=7.3cm 0cm 7.3cm 0cm,clip=true]{./images/new_gender_scatter.pdf}
% \end{subfigure}
% \begin{subfigure}[b]{0.33\textwidth}
% \includegraphics[width=\textwidth,trim=7.3cm 0cm 7.3cm 0cm,clip=true]{./images/new_relig_scatter.pdf}
% \end{subfigure}
% \caption{Examples of targets from the ``\emph{Profession}'', ``\emph{Gender}'', and ``\emph{Religion}'' categories from~\citet{nadeem2020stereoset} labeled by the regard measure. Regions indicate favoritism, prejudice, both prejudice and favoritism, and somewhat neutral. Higher negative regard percentages indicates prejudice-leaning, higher positive regard percentages indicates favoritism-leaning, and the region close to center representing somewhat fair region.}
% \label{fig:scatter_prof_sent}
% \end{figure*}



% \begin{figure*}[h]
% \centering
% \begin{subfigure}[b]{0.49\textwidth}
% \includegraphics[width=\textwidth,trim=0.3cm 8.6cm 0.2cm 8.3cm,clip=true]{./images/ex_race_2.pdf}
% \end{subfigure}
% \begin{subfigure}[b]{0.49\textwidth}
% \includegraphics[width=\textwidth,trim=0.3cm 0.1cm 0.2cm 0.3cm,clip=true]{./images/ex_gender_2.pdf}
% \end{subfigure}
% \begin{subfigure}[b]{0.49\textwidth}
% \includegraphics[width=\textwidth,trim=0.3cm 0.2cm 0.1cm 0.3cm,clip=true]{./images/ex_relig_2.pdf}
% \end{subfigure}
% \begin{subfigure}[b]{0.49\textwidth}
% \includegraphics[width=\textwidth,trim=0.3cm 0.2cm 0.2cm 0cm,clip=true]{./images/ex_prof_2.pdf}
% \end{subfigure}
% % \begin{subfigure}[b]{0.15\textwidth}
% % \includegraphics[width=\textwidth,trim=0cm 0cm 0cm 0cm,clip=true]{./images/conceptnet_bar_plot_sent_prof.pdf}
% % \end{subfigure}
% % \begin{subfigure}[b]{0.15\textwidth}
% % \includegraphics[width=\textwidth,trim=0cm 0cm 0cm 0cm,clip=true]{./images/conceptnet_bar_plot_regard_prof.pdf}
% % \end{subfigure}
% \caption{Four different representations from four categories each demonstrating a certain aspect of bias. In ``\emph{Origin}'' category, we can observe extreme overgeneralization toward ``\emph{british}'', in ``\emph{Gender}'' category both target groups are overgeneralized, in ``\emph{Religion}'' extreme prejudice toward ``\emph{muslim}'', and in ``\emph{Profession}'' extreme favoritism toward ``\emph{teacher}'' target group. Each case is accompanied with an example of negative and positive associations detected by sentiment.}
% \label{fig:sample_Bias_results}
% \end{figure*}




\para{Inter-target Disparity}
In addition to overgeneralized non-neutral views for each target group, we also study \textit{inter-target disparity} -- \textit{i.e.}, how different a target $t_j$ is perceived in the set $\mathbb{S}$ compared to other targets. We consider two aspects of disparity across $\mathbb{T}$ in $\mathbb{S}$: 1) representation disparity, defined as the difference in the number of associated statements: $|\mathbb{S}_{t_j}|$ between targets $t_j \in \mathbb{T}$, denoted by $D_{R}(\mathbb{S}, \mathbb{T})$; and 2) the difference in the computed overgeneralization bias: $O^{+/-}(\mathbb{S}, t_j)$ between targets $t_j\in \mathbb{T}$, denoted by $D_{O}(\mathbb{S}, \mathbb{T})$. Note that compared to $O^{+/-}(\mathbb{S}, t_j)$, $D_{O}(\mathbb{S}, \mathbb{T})$ is calculated over the full population of targets, thus measuring \emph{inter-target} disparity.
%Representation disparity measures if there is a severe under-representation for certain targets in the sentence set and overgeneralization disparity examines if some targets in general are perceived more positively than other targets, which are both indicators that the sentence set contains biases towards targets it depicts.
% \xiang{I would hope to have 1-2 sentences to explain the motivation and intuition of these two metrics; for example, is ``representation disparity" more about popularity or reporting bias?}
%We consider both of them as harmful disparities as it is both biased to under-represent certain targets and to make some groups more positively or negatively perceived than others.
For both aspects, we measure disparity using variance as follows:
\begin{equation} \label{eq:3}
      D_{R}(\mathbb{S}, \mathbb{T}) =\mathbb{E}[(|\mathbb{S}_{t_j}|-|\overline{\mathbb{S}_{t}}|)^2],
    %   \vspace{-1em}
\end{equation}
\begin{equation} \label{eq:4}
 D_{O}^{+/-}(\mathbb{S}, \mathbb{T}) =\mathbb{E}[(O^{+/-}(\mathbb{S}, t_j)-\overline{O^{+/-}(\mathbb{S}, t_j)})^2],
\end{equation}
where $|\overline{\mathbb{S}_{t}}|$ indicates the average number of statements for targets in $\mathbb{T}$ and $\overline{O^{+/-}(\mathbb{S}, t_j)}$ is the average overgeneralization bias for targets, ``+'' for favoritism and ``-'' for prejudice. The expectation $\mathbb{E}$ is taken over all targets $t_j \in \mathbb{T}$.




\subsection{Measuring Polarized Perceptions}\label{measure_quality}

%\xiang{specific implementation of the measurement of harms and the quality check/validation goes in this subsection.}
Prior work~\cite{sheng2019woman} demonstrated that sentiment and regard are effective measures of bias (polarized views toward a target group).
Although this is still an active area of research, for now, these are promising proxies that many works in ethical NLP also have used to measure bias (e.g.~\citet{sheng2019woman, li2020unqover, NEURIPS2020_1457c0d6, sheng-etal-2020-towards,dhamala2021bold}). However, we acknowledge that there still exist problems with these measures as proxies for measuring bias and acknowledge the existence of noisy labels using these measures as proxies. To put this into test and to show that these measures can still be reliable proxies despite the aforementioned problems, we perform studies both including human evaluators in the loop as well as comparison of these measures with a keyword-based approach in this section. 

In order to determine the polarization of perception associated to a statement toward a group, we apply sentiment and regard classifiers on the statement containing the target group and obtain the corresponding labels from each of the classifiers. We then categorize the statement into favoritism, prejudice, or neutral based on the positive, negative, or neutral labels obtained from each of the classifiers. 
%More experimental details on measuring polarized perceptions can be found in Section~\ref{sec:3} and Appendix Section D.

\para{Crowdsourcing Human Labels}
To validate the quality of these polarity proxies, we conduct crowdsourcing to solicit human labels on the statement polarity.
We asked Amazon Mechanical Turk workers to label provided knowledge from GenericsKB~\cite{bhakthavatsalam2020genericskb} and ConceptNet~\cite{speer2017conceptnet} with regards to favoritism, prejudice, and neutral toward a target group. 3,000 instances were labeled from ConceptNet and more than 1,500 from GenericsKB. The inter-annotator agreement in terms of Fleiss' kappa scores \cite{fleiss1971measuring} for this task was 0.5007 and 0.3827 for GenericsKB and ConceptNet respectively. 

\begin{table}[tb]
\scalebox{0.73}{
    \begin{tabular}{c | c c  | c c }
        \toprule
        CSKB & \multicolumn{2}{c}{GenericsKB} & \multicolumn{2}{c}{ConceptNet} \\
        \midrule
        Measure & Sentiment & Regard   & Sentiment & Regard \\
        \midrule
        Human Agreement &70.3\% & 60.9\% & 83.1\% & 75.4\%\\
        \bottomrule
    \end{tabular}}
    \caption{Agreement of sentiment and regard labels with human annotators in terms of accuracy.}
    \label{human_vs_re_sent_generics}
\end{table}


\begin{table}[tb]
\centering
\scalebox{0.72}{
\begin{tabular}{p{1.8cm}|ccc}
 \toprule
\textbf{CSKB}&\textbf{Method}&\textbf{Favoritism R/P/F1}&\textbf{Prejudice R/P/F1}\\[0.5pt]
 \midrule
 \parbox[t]{2mm}{\multirow{3}{*}{\shortstack[c]{GenericsKB}}}
&Regard&\textbf{0.551}/0.579/\textbf{0.565}&\textbf{0.809}/0.333/0.472 \\[0.5pt]
&Sentiment&0.441/0.622/0.516&0.432/\textbf{0.541}/\textbf{0.480}\\[0.5pt]
&Keyword&0.268/\textbf{0.643}/0.379&0.276/0.539/0.365\\[0.5pt]
 \midrule
  \parbox[t]{2mm}{\multirow{2}{*}{\shortstack[c]{ConceptNet}}}
  &Regard&\textbf{0.436}/0.383/0.408&\textbf{0.698}/0.342/\textbf{0.459}\\[0.5pt]
  &Sentiment&0.378/0.528/\textbf{0.440}&0.264/\textbf{0.531}/0.353\\[0.5pt]
 &Keyword&0.201/\textbf{0.556}/0.295&0.105/0.470/0.172\\[0.5pt]
  \bottomrule
\end{tabular}
}

\caption{Comparison of sentiment, regard, and baseline keyword-based approach in terms of favoritism and prejudice recall/precision/F1 scores.}
\label{recall}
\end{table}
\para{Alignment with Human Labels}
We compare human labels with those obtained from sentiment and regard classifiers to check the validity of these measures as proxies for overgeneralization. As shown in Table \ref{human_vs_re_sent_generics}, we found reasonable agreement in terms of accuracy for sentiment and regard with human labels. This was also confirmed in previous work \cite{sheng2019woman} in which sentiment and regard were shown to be good proxies to measure bias. 

\para{Comparison with Keyword-based Approach}
We also compare the sentiment and regard classifiers to a keyword-based baseline, in which we collect a list of biased words that could represent favoritism and prejudice from LIWC \cite{doi:10.1177/0261927X09351676} and Empath \cite{10.1145/2858036.2858535}. 
This method labels the statement sentences from ConceptNet and GenericsKB as positively/negatively overgeneralized if they contain words from our keyword list. As shown in Table \ref{recall}, this method has a significantly lower recall and overall F1 value in identifying favoritism and prejudice compared to sentiment and regard measures.
%we analyze the sufficiency and representativeness of triples associated to different demographic groups in ConceptNet. This problem can be observed as a data imbalance problem in knowledge graphs which can cause serious issues in many downstream tasks. Having a uniform and balanced knowledge coverage for different demographics is a critical principle as imbalance in existing knowledge can help systems to perform better for certain demographics that have more available and accessible knowledge in the knowledge graph. Imbalances in the existing knowledge can create disparities in performance of downstream tasks that rely on the commonsense knowledge and can create biases towards certain demographics. To quantify the data imbalance problem in ConceptNet, we reported the frequency of existing triples for different demographic groups and analyzed the differences amongst them.


% \begin{figure}[h]
% \includegraphics[width=0.45\textwidth,trim=0cm 0cm 0cm 0cm,clip=true]{./images/conceptnet_frequency.pdf}
% \caption{ConceptNet frequency analysis.}
% \label{fig:ConceptNet_freq}
% \end{figure}

